\label{section:classical}
In this section we will introduce the classical semi-online model introduced by \cite{karp}

\begin{definition}[Semi-Online Model]
	In this model the two anticliques of the bipartite graph $G=(U,V,E)$ are divided into offline and online agents.
	Offline agents $U$ are vertices that belong to one anticlique and are always visible.
	Online agents $V$ are the vertices from the other anticlique and arrive in the online fashion.
	When an online agent arrives edges incident to it are revealed. It has to be either discarded or matched to an offline agent and that match cannot change later.
	The adversary is oblivious
\end{definition}
For the purposes of minimizing repetition we state that properties of this model, mainly: $U$ being always revealed, the matching being instant and irrevocable, the adversary being oblivious are present in all subsequent model unless specified otherwise

This structure, where one anticlique is always known, will be retained by all models we will present. It is not a necessity, a Fully-Online model where both anticliques arrive in an online fashion exists, but it out of the scope of this work as it is known to be strictly harder than the Semi-Online model.


The adversary is oblivious as an adaptive adversary is trivial.
\begin{remark}
	No algorithm has a competitiveness asymptotically higher than $\frac{1}{2}$ against an adaptive adversary
\end{remark}
\begin{figure}[H]
	\centering
	\begin{tikzpicture}
	\begin{pgfonlayer}{nodelayer}
		\node [style=default] (0) at (-2, 7) {};
		\node [style=default] (1) at (2, 7) {};
		\node [style=default] (2) at (-2, 6) {};
		\node [style=default] (3) at (2, 6) {};
		\node [style=default] (4) at (-2, 5) {};
		\node [style=default] (5) at (2, 5) {};
		\node [style=default] (6) at (-2, 3) {};
		\node [style=default] (7) at (2, 3) {};
		\node [style=default] (8) at (-2, 2) {};
		\node [style=default] (9) at (2, 2) {};
		\node [style=nil] (10) at (-2, 4.25) {.};
		\node [style=nil] (11) at (-2, 4) {.};
		\node [style=nil] (12) at (-2, 3.75) {.};
		\node [style=nil] (13) at (2, 4.25) {.};
		\node [style=nil] (14) at (2, 4) {.};
		\node [style=nil] (15) at (2, 3.75) {.};
		\node [style=nil] (16) at (-2.5, 7) {$u_1$};
		\node [style=nil] (17) at (-2.5, 6) {$u_2$};
		\node [style=nil] (18) at (-2.5, 5) {$u_3$};
		\node [style=nil] (19) at (-2.5, 2) {$u_{2n}$};
	\end{pgfonlayer}
	\begin{pgfonlayer}{edgelayer}
		\draw (0) to (1);
		\draw (2) to (3);
		\draw (4) to (5);
		\draw (6) to (7);
		\draw (8) to (9);
		\draw (2) to (1);
		\draw (4) to (3);
		\draw (4) to (1);
		\draw (6) to (1);
		\draw (6) to (5);
		\draw (8) to (7);
		\draw (8) to (5);
		\draw (8) to (3);
		\draw [style=greenl] (6) to (3);
		\draw [style=greenl] (1) to (8);
	\end{pgfonlayer}
\end{tikzpicture}

\end{figure}
\begin{proof}
Begin with $2n$ offline agents $V$ and $k=0,X_0=\emptyset$.\\
Adversary repeats the following $n$ times: reveals an agent $v$ connected to $U\backslash X_k$. $k=k+1$, for some unmatched $u+k\in U; U_{k+1}=U_k\cup\{u_k\}$.\\
Then adversary reveals $n$ agents denoted $V'$ connected to already matched offline agents. The ultimate size of the matching is $n$.

In the final graph, all matched vertices $V'$ are connected to all online agents and all unmatched $v_i$ agents are connected to all $u_j$ such that $j\geq i$. Clearly, it is possible to match $v_i$ to $u_i$ and $V'$ to $U'$, forming a matching of size $2n$\\
\end{proof}
\begin{remark}
	In the semi-online model, any deterministic algorithm has a competitiveness no greater than $\frac{1}{2}$, there is a deterministic algorithm with competitiveness of $\frac{1}{2}$
\end{remark}
\begin{algorithm}
	\caption{Naive}
\begin{algorithmic}
	\item on Arrival of vertex $v$: match it with any offline agent if available
\end{algorithmic}
\end{algorithm}
\begin{proof}
	Trivially the competitiveness cannot be greater than $\frac{1}{2}$ as against a deterministic algorithm there is no difference between an adaptive and an oblivious adversary.

	To prove that there is a deterministic algorithm with competitiveness of $\frac{1}{2}$ consider the a naive algorithm.
	
	for every $(u,v)\in E(G)$ either $u$ or $v$ is matched. Therefore at least half of all vertices that could be matched are matched, therefore the competitiveness isn't lower than $\frac{1}{2}$
\end{proof}

The main discovery of \cite{karp} is the following algorithm

\begin{algorithm}
	\caption{Ranking}
	\label{alg:ranking}
	\begin{algorithmic}
		\item on Initialization: create a random permutation of the offline agents $\sigma$
		\item on Arrival of vertex $v$: match it with the first available offline neighbor according to $\sigma$
	\end{algorithmic}
\end{algorithm}

as well as the following theorems

\begin{theorem}
	\label{thm:rank-competitiveness}
	Ranking has asymptotic competitiveness of $1-\frac{1}{e}$
\end{theorem}
\begin{theorem}
	\label{thm:rank-optimal}
	No algorithm has asymptotic competitiveness higher than $1-\frac{1}{e}$
\end{theorem}

\section{Ranking Competitiveness}
	In this section we present a simple proof for \autoref{thm:rank-competitiveness} by \cite{obmSimple}. Many proofs of this theorem use linear programming, but this one uses no special tools

\begin{lemma}
	\label{rank-cmp-compete-perfect}
	Competitiveness is determined by graphs with a perfect matching
\end{lemma}
\begin{proof}
	Let $G$ be a bipartite graph without a perfect matching, with orderings of offline agents $\sigma$ and of online agents $\pi$. Consider $G'$ which is a copy of $G$ without a vertex $x$ and orderings $\pi'$ induced by $V\backslash\{x\}$.

	Observe that $m=\mathrm{Ranking}(G,\pi,\sigma)$ being the matching found by the Ranking algorithm on $G$ with ranking $\pi$ and arrival order $\sigma$ is either identical to $m'=\mathrm{Ranking}(G',\pi',\sigma')$ or has one more vertex pair. This is apparent since either:
	\begin{enumerate}
		\item $x$ isn't in $m$,
		\item $x' = M(x)$ isn't matched in $m'$, or
		\item $x'$ is matched to $y$ in $m'$
	\end{enumerate}
	In the case of 3 there is an alternating path of edges in $m'$ and those that are in $m$ but not in $m'$. When this path ends in the same partition as $x$ then $m$ is larger than $m'$ by 1

	This also leads to the conclusion that if the respective matching sizes are different, they differ by  a single alternating path.

	Now consider the maximal matching of $G$ and the effect the removal of the vertices outside of it will have on the competitiveness.
	The size of the maximal matching will (of course) not change, while the Ranking's solution decreases by at most 1 with each removal, implying that this process of reducing $G$ to $G_p$ with a perfect matching didn't increase the competitiveness, meaning that analysis of graphs without perfect matchings is redundant.
\end{proof}

In the following section graphs are assumed to have perfect matchings. Let $M^\star$ be that perfect matching.

\begin{lemma}
	\label{rank-cmp-unmatched-simple}
	Let $v$ online vertex,$u=m^\star(v)$, $\sigma$ permutation.\\
	If $u$ isn't matched by $\mathrm{Ranking}(\sigma)$ then $v$ is matched by $\mathrm{Ranking}(\sigma)$ to $u'$ with a lower rank
\end{lemma}

\begin{proof}
	Trivial since $u,v$ are connected, therefore for $v$ to remain unmatched $u$ must've been matched to another vertex, which must've in turn had a lower rank than $v$
\end{proof}

\begin{lemma}
	\label{rank-cmp-unmatched-complicated}
	Let $v$ online vertex,$u=m^\star(v)$, $\sigma$ permutation, $\sigma_i$ permutation obtained by removing $v$ and inserting it into $sigma$ at the $i$-th place.\\
	If $u$ isn't matched by $\mathrm{Ranking}(\sigma)$ then for every $i$ $\mathrm{Ranking}(\sigma_i)$ matches $v$ with a vertex $u_i$ such that $\sigma_i(v_i)\leq\sigma(v)=t$
\end{lemma}

\begin{proof}
	$m=\mathrm{Ranking}(\sigma),m_i=\mathrm{Ranking}(\sigma_i)$.\\ Symmetrically to \autoref{rank-cmp-compete-perfect} $m,m_i$ are either identical or differ by a single alternating path. By repeatedly applying the logic from the proof of \autoref{rank-cmp-unmatched-simple} see that offline vertices in that path are of increasing rank in $\sigma_i$. Therefore $v_i=m_i(u),v'=m(u),\sigma_i(v_i)\leq\sigma_i(v')$\\
	Using \autoref{rank-cmp-unmatched-simple} we get $\sigma(v')<t$ and since $|\sigma_i(v')-\sigma(v')|\leq 1$ we receive $\sigma_i(v_i)<1+t$ equivalent to $\sigma_i(v_i)\leq t$ on integers
\end{proof}

The final lemma required to complete the proof uses this result directly.

\begin{lemma}
	\label{rank-cmp-probability}
	For random $\sigma$, the probability that the offline vertex of rank $t$ is matched by $\mathrm{Ranking}(\sigma)$ denoted $x_t$ fulfills the following: $1-x_t\leq\frac{1}{n}\sum\limits_{1\leq s\leq t}x_s$
\end{lemma}

\begin{proof}
	Let $\sigma$ be a random permutation of offline vertices. Define $\sigma'$ with the following random process: choose an offline vertex randomly with equal probability, take it out of $\sigma$ and insert it into the $t$-th rank. Let $u$ be the vertex of rank $t,u=m^*(u),R_{t-1}$\\ online vertices matched by the algorithm to offline vertices of rank $\leq t-1$\\
	probability that $v$ is not matched is $1-x_t$. Expcted cardinality of $R_{t-1}$ is $\sum\limits_{1\leq s\leq t-1}x_s$.\\
	Applying \autoref{rank-cmp-unmatched-complicated} with $\sigma:=\sigma',i$ chosen such that $\sigma'_i=\sigma$\ : if $v$ isn't matched by $\mathrm{Ranking}(\sigma')$ then $v$ is matched by $\mathrm{Ranking}(\sigma)$ to a vertex $u'$ such that $\sigma(v)\leq t$ implying $v\in R_t$\\
	Conditional on $\sigma$, $v\in R_t$ has a probability of $\frac{|R_t|}{n}$ since they are independent.\\
	We now have two random events $E,F$ that fulfill the condition that $E\implies F$. It is obvious to see that $\Pr[E]\geq\Pr[F]$\\
	The lemma follows
	
\end{proof}

With \autoref{rank-cmp-probability} the Theorem can be proved directly.\\
$G$ has a perfect matching implies that copetitiveness is $\inf \frac{1}{n} \sum\limits_{1\leq s\leq n}x_s$.\\
Let $S_t=\sum\limits_{1\leq s\leq t}x_s$. \autoref{rank-cmp-probability} can be rewritten as $S_t(1+\frac{1}{n})\geq 1 + S_{t-1}$ and the competitiveness ratio as $\frac{1}{n}S_n$.\\
Infimum occurs therefore when all inequalities are tight, resulting in $S_t=\sum\limits_{1\leq s\leq t}(1-\frac{1}{n+1})^s$ for all $t$, implying that competitiveness is at least $\frac{1}{n}\sum\limits_{1\leq s\leq n}(1-\frac{1}{n+1})^s=1-(1-\frac{1}{n+1})^n\xrightarrow[n\to\infty]{}1-\frac{1}{e}$

Thus the theorem is proven
\section{Ranking Optimality}
	In this section we present a proof for Ranking Optimality (theorem 2) by \cite{fallback}